% from ega2-1-fr.tex
\oldpage[II]{170}
qui correspond au diagramme
\[
  \xymatrix{
    \mathcal{S}\ar[r]^{\varphi}\ar[d] & \mathcal{S}'\ar[d]\\
    \mathcal{O}_Y\ar[r] & \mathcal{O}_{Y'}
  }
\]
où les flèches verticales sont les homomorphismes d'augmentation, et dont la
commutativité résulte de l'hypothèse que $\varphi$ est supposé être un
homomorphisme d'Algèbres \emph{graduées}.
\end{env}

\begin{proposition}[8.5.2]
\label{II.8.5.2-fr}
Si $E$ (resp. $E'$) est le cône épointé affine défini par $\mathcal{S}$
(resp. $\mathcal{S}'$), on a $\Phi^{-1}(E)\subset E'$~; si en outre
$\operatorname{Proj}(\varphi):G(\varphi)\to\operatorname{Proj}(\mathcal{S})$
est partout défini (autrement dit si
$G(\varphi)=\operatorname{Proj}(\mathcal{S}')$), alors on a
$\Phi^{-1}(E)=E'$, et réciproquement.
\end{proposition}

\begin{proof}
La première assertion résulte de la commutativité de
(\hyperref[II.8.5.1.3-fr]{8.5.1.3}). Pour démontrer la seconde, on peut se
borner au cas où $Y=\operatorname{Spec}(A)$ et
$Y'=\operatorname{Spec}(A')$ sont affines,
$\mathcal{S}=\widetilde{S}$, $\mathcal{S}'=\widetilde{S'}$. Pour tout $f$
homogène dans $S_+$, si on pose $f'=\varphi(f)$, on a
$\Phi^{-1}(C_f)=C'_{f'}$ (\hyperref[I.2.2.4.1-fr]{I, 2.2.4.1})~; dire que
$G(\varphi)=\operatorname{Proj}(S')$ signifie que dans $S'_+$ la racine de
l'idéal engendré par les $f'=\varphi(f)$ est $S'_+$ lui-même
(\hyperref[II.2.8.1-fr]{2.8.1}) et (\hyperref[II.2.3.14-fr]{2.3.14}), et cela
équivaut encore à dire que les $C'_{f'}$ recouvrent $E'$
(\hyperref[II.8.3.4.4-fr]{8.3.4.4}).
\end{proof}

\begin{env}[8.5.3]
\label{II.8.5.3-fr}
Le $q$-morphisme $\varphi$ se prolonge canoniquement en un $q$-morphisme
d'Algèbres graduées
\begin{equation}
\label{II.8.5.3.1-fr}
  \widehat{\varphi}:\widehat{\mathcal{S}}
  \to\widehat{\mathcal{S}'}
\tag{8.5.3.1}
\end{equation}
en posant $\widehat{\varphi}(z)=z$. On en déduit par suite un morphisme
\[
  \widehat{\Phi}=\operatorname{Proj}(\widehat{\varphi}):
  G(\widehat{\varphi})\to\widehat{C}
  =\operatorname{Proj}(\widehat{\mathcal{S}})
\]
tel que le diagramme
\[
  \xymatrix{
    G(\widehat{\varphi})\ar[r]^{\widehat{\Phi}}\ar[d] &
    \widehat{C}\ar[d]\\
    Y'\ar[r]_{q} & Y
  }
\]
soit commutatif (\hyperref[II.3.5.6-fr]{3.5.6}). Il résulte aussitôt des
définitions que si on désigne par $i:C\to\widehat{C}$ et
$i':C'\to\widehat{C'}$ les immersions ouvertes canoniques
(\hyperref[II.8.3.2-fr]{8.3.2}), on a
$i'(C')\subset G(\widehat{\varphi})$ et le diagramme
\begin{equation}
\label{II.8.5.3.2-fr}
  \xymatrix{
    C'\ar[r]^{\Phi}\ar[d]_{i'} & C\ar[d]^{i}\\
    G(\widehat{\varphi})\ar[r]_{\widehat{\Phi}} & \widehat{C}
  }
\tag{8.5.3.2}
\end{equation}
est commutatif. Enfin, si on pose
$X=\operatorname{Proj}(\mathcal{S})$, $X'=\operatorname{Proj}(\mathcal{S}')$,
et si $j:X\to\widehat{C}$, $j':X'\to\widehat{C'}$ sont les immersions
fermées canoniques (\hyperref[II.8.3.2-fr]{8.3.2}), il résulte encore des
définitions de ces immersions que l'on a
$j'(G(\varphi))\subset G(\widehat{\varphi})$ et que le diagramme

